% ============================================================================
% ESQUEMA PROPOSTO (banca, sugestão — não mergeado): a geometria da LCE
% (Seção sec:metodo-metricas). A curva é ILUSTRATIVA e o desenho diz isso
% no título interno: nenhum dado medido, nenhum número além dos símbolos da
% própria equação — é a definição da métrica desenhada, não um resultado
% (compatível com o princípio V). Por decisão do autor (2026-08-23), a
% análise estatística e a precedência da ALC ficam SÓ na prosa.
%
% Uso no Cap. 3 (dentro de um ambiente figure):
%   % ============================================================================
% ESQUEMA PROPOSTO (banca, sugestão — não mergeado): a geometria da LCE
% (Seção sec:metodo-metricas). A curva é ILUSTRATIVA e o desenho diz isso
% no título interno: nenhum dado medido, nenhum número além dos símbolos da
% própria equação — é a definição da métrica desenhada, não um resultado
% (compatível com o princípio V). Por decisão do autor (2026-08-23), a
% análise estatística e a precedência da ALC ficam SÓ na prosa.
%
% Uso no Cap. 3 (dentro de um ambiente figure):
%   % ============================================================================
% ESQUEMA PROPOSTO (banca, sugestão — não mergeado): a geometria da LCE
% (Seção sec:metodo-metricas). A curva é ILUSTRATIVA e o desenho diz isso
% no título interno: nenhum dado medido, nenhum número além dos símbolos da
% própria equação — é a definição da métrica desenhada, não um resultado
% (compatível com o princípio V). Por decisão do autor (2026-08-23), a
% análise estatística e a precedência da ALC ficam SÓ na prosa.
%
% Uso no Cap. 3 (dentro de um ambiente figure):
%   % ============================================================================
% ESQUEMA PROPOSTO (banca, sugestão — não mergeado): a geometria da LCE
% (Seção sec:metodo-metricas). A curva é ILUSTRATIVA e o desenho diz isso
% no título interno: nenhum dado medido, nenhum número além dos símbolos da
% própria equação — é a definição da métrica desenhada, não um resultado
% (compatível com o princípio V). Por decisão do autor (2026-08-23), a
% análise estatística e a precedência da ALC ficam SÓ na prosa.
%
% Uso no Cap. 3 (dentro de um ambiente figure):
%   \input{3-metodo/esquemas-propostos/esq-lce}
%   A citação Guyon2011ALC (precedente ALC) permanece na prosa/legenda.
% Uso standalone (pré-visualização):
%   \documentclass[tikz,border=6pt]{standalone}
%   \usetikzlibrary{arrows.meta, positioning}
%   \begin{document}\input{esq-lce.tex}\end{document}
%   (no modo standalone, troque as \ref{...} por texto fixo ou defina-as)
% Requer apenas arrows.meta e positioning (já em packages.tex). Legível em
% P&B: área sob a curva real = cinza escuro; retângulo ideal = cinza claro;
% pontos observados marcados (é sobre eles que Simpson integra).
% ============================================================================
\begin{tikzpicture}[
    x=1cm, y=1cm,
    caixa/.style={draw, rounded corners=2pt, align=center, text width=42mm,
                  minimum height=11mm, inner sep=3pt, font=\footnotesize},
    seta/.style={-{Stealth[length=2.4mm]}, thick},
    rot/.style={font=\scriptsize, align=center, inner sep=2pt}
]
    % ---------------- painel esquerdo: a geometria -------------------------
    \node[rot, anchor=west] at (-0.4,5.2)
        {\textbf{A geometria da LCE} (curva ilustrativa, sem dados medidos)};

    % áreas (desenhadas antes dos eixos)
    \fill[gray!8]  (0.8,0) rectangle (6.4,3.8);
    \fill[gray!25] (0.8,1.0) -- (1.5,2.2) -- (2.4,2.9) -- (3.4,3.3)
                   -- (4.6,3.55) -- (6.4,3.7) -- (6.4,0) -- (0.8,0) -- cycle;
    \draw[thin] (0.8,0) rectangle (6.4,3.8);

    % curva de aprendizado com os pontos observados
    \draw[thick] (0.8,1.0) -- (1.5,2.2) -- (2.4,2.9) -- (3.4,3.3)
                 -- (4.6,3.55) -- (6.4,3.7);
    \foreach \p in {(0.8,1.0),(1.5,2.2),(2.4,2.9),(3.4,3.3),(4.6,3.55),(6.4,3.7)}
        \fill \p circle (1.6pt);

    % linha do baseline (rótulo à direita, longe do ponteiro da AUC_ideal)
    \draw[dashed] (0,3.8) -- (7.0,3.8);
    \node[rot, anchor=south east] at (6.95,3.86)
        {\mbox{$\mathrm{Perf}_{\mathrm{baseline}}$} (supervisão completa)};

    % eixos
    \draw[seta] (-0.15,0) -- (7.35,0);
    \draw[seta] (0,-0.15) -- (0,4.6);
    \node[rot, anchor=south] at (0.3,4.6) {desempenho};
    \node[rot, anchor=north] at (3.6,-0.62) {rótulos consumidos ($|L_t|$)};
    \draw (0.8,0) -- (0.8,-0.12);
    \draw (6.4,0) -- (6.4,-0.12);
    \node[rot, anchor=north] at (0.8,-0.2) {$L_{\mathrm{ideal},0}$};
    \node[rot, anchor=north] at (6.4,-0.2) {$L_{\mathrm{final}}$};

    % rótulos das áreas: a fatia clara é estreita demais para texto dentro,
    % então a AUC_ideal é apontada de fora (achado da inspeção visual)
    \node[rot] at (4.5,1.7) {$\mathrm{AUC}_{\mathrm{real}}$};
    \node[rot, anchor=west] (auci) at (0.7,4.32)
        {$\mathrm{AUC}_{\mathrm{ideal}}$: o retângulo inteiro};
    \draw[seta, dashed] (1.6,4.15) -- (1.4,3.5);

    % ---------------- painel direito: fórmula e leitura --------------------
    \node[caixa, text width=58mm] (formula) at (10.6,3.5)
        {$\displaystyle \mathrm{LCE}=\frac{\mathrm{AUC}_{\mathrm{real}}}{\mathrm{AUC}_{\mathrm{ideal}}}$\\[3pt]
         \mbox{$\mathrm{AUC}_{\mathrm{ideal}}=\mathrm{Perf}_{\mathrm{baseline}}\times(L_{\mathrm{final}}-L_{\mathrm{ideal},0})$}\\[2pt]
         (Equação~\ref{eq:metodo-lce}; Apêndice~\ref{ap:lce})};
    \node[rot, text width=52mm, anchor=north, align=left] at (10.6,2.2)
        {$\mathrm{LCE}\approx 1$: a estratégia aproxima-se\\
         rapidamente do desempenho de\\
         referência com poucos rótulos.\\
         Integração numérica: regra de Simpson\\
         sobre os pontos observados (marcados).};
\end{tikzpicture}

%   A citação Guyon2011ALC (precedente ALC) permanece na prosa/legenda.
% Uso standalone (pré-visualização):
%   \documentclass[tikz,border=6pt]{standalone}
%   \usetikzlibrary{arrows.meta, positioning}
%   \begin{document}% ============================================================================
% ESQUEMA PROPOSTO (banca, sugestão — não mergeado): a geometria da LCE
% (Seção sec:metodo-metricas). A curva é ILUSTRATIVA e o desenho diz isso
% no título interno: nenhum dado medido, nenhum número além dos símbolos da
% própria equação — é a definição da métrica desenhada, não um resultado
% (compatível com o princípio V). Por decisão do autor (2026-08-23), a
% análise estatística e a precedência da ALC ficam SÓ na prosa.
%
% Uso no Cap. 3 (dentro de um ambiente figure):
%   \input{3-metodo/esquemas-propostos/esq-lce}
%   A citação Guyon2011ALC (precedente ALC) permanece na prosa/legenda.
% Uso standalone (pré-visualização):
%   \documentclass[tikz,border=6pt]{standalone}
%   \usetikzlibrary{arrows.meta, positioning}
%   \begin{document}\input{esq-lce.tex}\end{document}
%   (no modo standalone, troque as \ref{...} por texto fixo ou defina-as)
% Requer apenas arrows.meta e positioning (já em packages.tex). Legível em
% P&B: área sob a curva real = cinza escuro; retângulo ideal = cinza claro;
% pontos observados marcados (é sobre eles que Simpson integra).
% ============================================================================
\begin{tikzpicture}[
    x=1cm, y=1cm,
    caixa/.style={draw, rounded corners=2pt, align=center, text width=42mm,
                  minimum height=11mm, inner sep=3pt, font=\footnotesize},
    seta/.style={-{Stealth[length=2.4mm]}, thick},
    rot/.style={font=\scriptsize, align=center, inner sep=2pt}
]
    % ---------------- painel esquerdo: a geometria -------------------------
    \node[rot, anchor=west] at (-0.4,5.2)
        {\textbf{A geometria da LCE} (curva ilustrativa, sem dados medidos)};

    % áreas (desenhadas antes dos eixos)
    \fill[gray!8]  (0.8,0) rectangle (6.4,3.8);
    \fill[gray!25] (0.8,1.0) -- (1.5,2.2) -- (2.4,2.9) -- (3.4,3.3)
                   -- (4.6,3.55) -- (6.4,3.7) -- (6.4,0) -- (0.8,0) -- cycle;
    \draw[thin] (0.8,0) rectangle (6.4,3.8);

    % curva de aprendizado com os pontos observados
    \draw[thick] (0.8,1.0) -- (1.5,2.2) -- (2.4,2.9) -- (3.4,3.3)
                 -- (4.6,3.55) -- (6.4,3.7);
    \foreach \p in {(0.8,1.0),(1.5,2.2),(2.4,2.9),(3.4,3.3),(4.6,3.55),(6.4,3.7)}
        \fill \p circle (1.6pt);

    % linha do baseline (rótulo à direita, longe do ponteiro da AUC_ideal)
    \draw[dashed] (0,3.8) -- (7.0,3.8);
    \node[rot, anchor=south east] at (6.95,3.86)
        {\mbox{$\mathrm{Perf}_{\mathrm{baseline}}$} (supervisão completa)};

    % eixos
    \draw[seta] (-0.15,0) -- (7.35,0);
    \draw[seta] (0,-0.15) -- (0,4.6);
    \node[rot, anchor=south] at (0.3,4.6) {desempenho};
    \node[rot, anchor=north] at (3.6,-0.62) {rótulos consumidos ($|L_t|$)};
    \draw (0.8,0) -- (0.8,-0.12);
    \draw (6.4,0) -- (6.4,-0.12);
    \node[rot, anchor=north] at (0.8,-0.2) {$L_{\mathrm{ideal},0}$};
    \node[rot, anchor=north] at (6.4,-0.2) {$L_{\mathrm{final}}$};

    % rótulos das áreas: a fatia clara é estreita demais para texto dentro,
    % então a AUC_ideal é apontada de fora (achado da inspeção visual)
    \node[rot] at (4.5,1.7) {$\mathrm{AUC}_{\mathrm{real}}$};
    \node[rot, anchor=west] (auci) at (0.7,4.32)
        {$\mathrm{AUC}_{\mathrm{ideal}}$: o retângulo inteiro};
    \draw[seta, dashed] (1.6,4.15) -- (1.4,3.5);

    % ---------------- painel direito: fórmula e leitura --------------------
    \node[caixa, text width=58mm] (formula) at (10.6,3.5)
        {$\displaystyle \mathrm{LCE}=\frac{\mathrm{AUC}_{\mathrm{real}}}{\mathrm{AUC}_{\mathrm{ideal}}}$\\[3pt]
         \mbox{$\mathrm{AUC}_{\mathrm{ideal}}=\mathrm{Perf}_{\mathrm{baseline}}\times(L_{\mathrm{final}}-L_{\mathrm{ideal},0})$}\\[2pt]
         (Equação~\ref{eq:metodo-lce}; Apêndice~\ref{ap:lce})};
    \node[rot, text width=52mm, anchor=north, align=left] at (10.6,2.2)
        {$\mathrm{LCE}\approx 1$: a estratégia aproxima-se\\
         rapidamente do desempenho de\\
         referência com poucos rótulos.\\
         Integração numérica: regra de Simpson\\
         sobre os pontos observados (marcados).};
\end{tikzpicture}
\end{document}
%   (no modo standalone, troque as \ref{...} por texto fixo ou defina-as)
% Requer apenas arrows.meta e positioning (já em packages.tex). Legível em
% P&B: área sob a curva real = cinza escuro; retângulo ideal = cinza claro;
% pontos observados marcados (é sobre eles que Simpson integra).
% ============================================================================
\begin{tikzpicture}[
    x=1cm, y=1cm,
    caixa/.style={draw, rounded corners=2pt, align=center, text width=42mm,
                  minimum height=11mm, inner sep=3pt, font=\footnotesize},
    seta/.style={-{Stealth[length=2.4mm]}, thick},
    rot/.style={font=\scriptsize, align=center, inner sep=2pt}
]
    % ---------------- painel esquerdo: a geometria -------------------------
    \node[rot, anchor=west] at (-0.4,5.2)
        {\textbf{A geometria da LCE} (curva ilustrativa, sem dados medidos)};

    % áreas (desenhadas antes dos eixos)
    \fill[gray!8]  (0.8,0) rectangle (6.4,3.8);
    \fill[gray!25] (0.8,1.0) -- (1.5,2.2) -- (2.4,2.9) -- (3.4,3.3)
                   -- (4.6,3.55) -- (6.4,3.7) -- (6.4,0) -- (0.8,0) -- cycle;
    \draw[thin] (0.8,0) rectangle (6.4,3.8);

    % curva de aprendizado com os pontos observados
    \draw[thick] (0.8,1.0) -- (1.5,2.2) -- (2.4,2.9) -- (3.4,3.3)
                 -- (4.6,3.55) -- (6.4,3.7);
    \foreach \p in {(0.8,1.0),(1.5,2.2),(2.4,2.9),(3.4,3.3),(4.6,3.55),(6.4,3.7)}
        \fill \p circle (1.6pt);

    % linha do baseline (rótulo à direita, longe do ponteiro da AUC_ideal)
    \draw[dashed] (0,3.8) -- (7.0,3.8);
    \node[rot, anchor=south east] at (6.95,3.86)
        {\mbox{$\mathrm{Perf}_{\mathrm{baseline}}$} (supervisão completa)};

    % eixos
    \draw[seta] (-0.15,0) -- (7.35,0);
    \draw[seta] (0,-0.15) -- (0,4.6);
    \node[rot, anchor=south] at (0.3,4.6) {desempenho};
    \node[rot, anchor=north] at (3.6,-0.62) {rótulos consumidos ($|L_t|$)};
    \draw (0.8,0) -- (0.8,-0.12);
    \draw (6.4,0) -- (6.4,-0.12);
    \node[rot, anchor=north] at (0.8,-0.2) {$L_{\mathrm{ideal},0}$};
    \node[rot, anchor=north] at (6.4,-0.2) {$L_{\mathrm{final}}$};

    % rótulos das áreas: a fatia clara é estreita demais para texto dentro,
    % então a AUC_ideal é apontada de fora (achado da inspeção visual)
    \node[rot] at (4.5,1.7) {$\mathrm{AUC}_{\mathrm{real}}$};
    \node[rot, anchor=west] (auci) at (0.7,4.32)
        {$\mathrm{AUC}_{\mathrm{ideal}}$: o retângulo inteiro};
    \draw[seta, dashed] (1.6,4.15) -- (1.4,3.5);

    % ---------------- painel direito: fórmula e leitura --------------------
    \node[caixa, text width=58mm] (formula) at (10.6,3.5)
        {$\displaystyle \mathrm{LCE}=\frac{\mathrm{AUC}_{\mathrm{real}}}{\mathrm{AUC}_{\mathrm{ideal}}}$\\[3pt]
         \mbox{$\mathrm{AUC}_{\mathrm{ideal}}=\mathrm{Perf}_{\mathrm{baseline}}\times(L_{\mathrm{final}}-L_{\mathrm{ideal},0})$}\\[2pt]
         (Equação~\ref{eq:metodo-lce}; Apêndice~\ref{ap:lce})};
    \node[rot, text width=52mm, anchor=north, align=left] at (10.6,2.2)
        {$\mathrm{LCE}\approx 1$: a estratégia aproxima-se\\
         rapidamente do desempenho de\\
         referência com poucos rótulos.\\
         Integração numérica: regra de Simpson\\
         sobre os pontos observados (marcados).};
\end{tikzpicture}

%   A citação Guyon2011ALC (precedente ALC) permanece na prosa/legenda.
% Uso standalone (pré-visualização):
%   \documentclass[tikz,border=6pt]{standalone}
%   \usetikzlibrary{arrows.meta, positioning}
%   \begin{document}% ============================================================================
% ESQUEMA PROPOSTO (banca, sugestão — não mergeado): a geometria da LCE
% (Seção sec:metodo-metricas). A curva é ILUSTRATIVA e o desenho diz isso
% no título interno: nenhum dado medido, nenhum número além dos símbolos da
% própria equação — é a definição da métrica desenhada, não um resultado
% (compatível com o princípio V). Por decisão do autor (2026-08-23), a
% análise estatística e a precedência da ALC ficam SÓ na prosa.
%
% Uso no Cap. 3 (dentro de um ambiente figure):
%   % ============================================================================
% ESQUEMA PROPOSTO (banca, sugestão — não mergeado): a geometria da LCE
% (Seção sec:metodo-metricas). A curva é ILUSTRATIVA e o desenho diz isso
% no título interno: nenhum dado medido, nenhum número além dos símbolos da
% própria equação — é a definição da métrica desenhada, não um resultado
% (compatível com o princípio V). Por decisão do autor (2026-08-23), a
% análise estatística e a precedência da ALC ficam SÓ na prosa.
%
% Uso no Cap. 3 (dentro de um ambiente figure):
%   \input{3-metodo/esquemas-propostos/esq-lce}
%   A citação Guyon2011ALC (precedente ALC) permanece na prosa/legenda.
% Uso standalone (pré-visualização):
%   \documentclass[tikz,border=6pt]{standalone}
%   \usetikzlibrary{arrows.meta, positioning}
%   \begin{document}\input{esq-lce.tex}\end{document}
%   (no modo standalone, troque as \ref{...} por texto fixo ou defina-as)
% Requer apenas arrows.meta e positioning (já em packages.tex). Legível em
% P&B: área sob a curva real = cinza escuro; retângulo ideal = cinza claro;
% pontos observados marcados (é sobre eles que Simpson integra).
% ============================================================================
\begin{tikzpicture}[
    x=1cm, y=1cm,
    caixa/.style={draw, rounded corners=2pt, align=center, text width=42mm,
                  minimum height=11mm, inner sep=3pt, font=\footnotesize},
    seta/.style={-{Stealth[length=2.4mm]}, thick},
    rot/.style={font=\scriptsize, align=center, inner sep=2pt}
]
    % ---------------- painel esquerdo: a geometria -------------------------
    \node[rot, anchor=west] at (-0.4,5.2)
        {\textbf{A geometria da LCE} (curva ilustrativa, sem dados medidos)};

    % áreas (desenhadas antes dos eixos)
    \fill[gray!8]  (0.8,0) rectangle (6.4,3.8);
    \fill[gray!25] (0.8,1.0) -- (1.5,2.2) -- (2.4,2.9) -- (3.4,3.3)
                   -- (4.6,3.55) -- (6.4,3.7) -- (6.4,0) -- (0.8,0) -- cycle;
    \draw[thin] (0.8,0) rectangle (6.4,3.8);

    % curva de aprendizado com os pontos observados
    \draw[thick] (0.8,1.0) -- (1.5,2.2) -- (2.4,2.9) -- (3.4,3.3)
                 -- (4.6,3.55) -- (6.4,3.7);
    \foreach \p in {(0.8,1.0),(1.5,2.2),(2.4,2.9),(3.4,3.3),(4.6,3.55),(6.4,3.7)}
        \fill \p circle (1.6pt);

    % linha do baseline (rótulo à direita, longe do ponteiro da AUC_ideal)
    \draw[dashed] (0,3.8) -- (7.0,3.8);
    \node[rot, anchor=south east] at (6.95,3.86)
        {\mbox{$\mathrm{Perf}_{\mathrm{baseline}}$} (supervisão completa)};

    % eixos
    \draw[seta] (-0.15,0) -- (7.35,0);
    \draw[seta] (0,-0.15) -- (0,4.6);
    \node[rot, anchor=south] at (0.3,4.6) {desempenho};
    \node[rot, anchor=north] at (3.6,-0.62) {rótulos consumidos ($|L_t|$)};
    \draw (0.8,0) -- (0.8,-0.12);
    \draw (6.4,0) -- (6.4,-0.12);
    \node[rot, anchor=north] at (0.8,-0.2) {$L_{\mathrm{ideal},0}$};
    \node[rot, anchor=north] at (6.4,-0.2) {$L_{\mathrm{final}}$};

    % rótulos das áreas: a fatia clara é estreita demais para texto dentro,
    % então a AUC_ideal é apontada de fora (achado da inspeção visual)
    \node[rot] at (4.5,1.7) {$\mathrm{AUC}_{\mathrm{real}}$};
    \node[rot, anchor=west] (auci) at (0.7,4.32)
        {$\mathrm{AUC}_{\mathrm{ideal}}$: o retângulo inteiro};
    \draw[seta, dashed] (1.6,4.15) -- (1.4,3.5);

    % ---------------- painel direito: fórmula e leitura --------------------
    \node[caixa, text width=58mm] (formula) at (10.6,3.5)
        {$\displaystyle \mathrm{LCE}=\frac{\mathrm{AUC}_{\mathrm{real}}}{\mathrm{AUC}_{\mathrm{ideal}}}$\\[3pt]
         \mbox{$\mathrm{AUC}_{\mathrm{ideal}}=\mathrm{Perf}_{\mathrm{baseline}}\times(L_{\mathrm{final}}-L_{\mathrm{ideal},0})$}\\[2pt]
         (Equação~\ref{eq:metodo-lce}; Apêndice~\ref{ap:lce})};
    \node[rot, text width=52mm, anchor=north, align=left] at (10.6,2.2)
        {$\mathrm{LCE}\approx 1$: a estratégia aproxima-se\\
         rapidamente do desempenho de\\
         referência com poucos rótulos.\\
         Integração numérica: regra de Simpson\\
         sobre os pontos observados (marcados).};
\end{tikzpicture}

%   A citação Guyon2011ALC (precedente ALC) permanece na prosa/legenda.
% Uso standalone (pré-visualização):
%   \documentclass[tikz,border=6pt]{standalone}
%   \usetikzlibrary{arrows.meta, positioning}
%   \begin{document}% ============================================================================
% ESQUEMA PROPOSTO (banca, sugestão — não mergeado): a geometria da LCE
% (Seção sec:metodo-metricas). A curva é ILUSTRATIVA e o desenho diz isso
% no título interno: nenhum dado medido, nenhum número além dos símbolos da
% própria equação — é a definição da métrica desenhada, não um resultado
% (compatível com o princípio V). Por decisão do autor (2026-08-23), a
% análise estatística e a precedência da ALC ficam SÓ na prosa.
%
% Uso no Cap. 3 (dentro de um ambiente figure):
%   \input{3-metodo/esquemas-propostos/esq-lce}
%   A citação Guyon2011ALC (precedente ALC) permanece na prosa/legenda.
% Uso standalone (pré-visualização):
%   \documentclass[tikz,border=6pt]{standalone}
%   \usetikzlibrary{arrows.meta, positioning}
%   \begin{document}\input{esq-lce.tex}\end{document}
%   (no modo standalone, troque as \ref{...} por texto fixo ou defina-as)
% Requer apenas arrows.meta e positioning (já em packages.tex). Legível em
% P&B: área sob a curva real = cinza escuro; retângulo ideal = cinza claro;
% pontos observados marcados (é sobre eles que Simpson integra).
% ============================================================================
\begin{tikzpicture}[
    x=1cm, y=1cm,
    caixa/.style={draw, rounded corners=2pt, align=center, text width=42mm,
                  minimum height=11mm, inner sep=3pt, font=\footnotesize},
    seta/.style={-{Stealth[length=2.4mm]}, thick},
    rot/.style={font=\scriptsize, align=center, inner sep=2pt}
]
    % ---------------- painel esquerdo: a geometria -------------------------
    \node[rot, anchor=west] at (-0.4,5.2)
        {\textbf{A geometria da LCE} (curva ilustrativa, sem dados medidos)};

    % áreas (desenhadas antes dos eixos)
    \fill[gray!8]  (0.8,0) rectangle (6.4,3.8);
    \fill[gray!25] (0.8,1.0) -- (1.5,2.2) -- (2.4,2.9) -- (3.4,3.3)
                   -- (4.6,3.55) -- (6.4,3.7) -- (6.4,0) -- (0.8,0) -- cycle;
    \draw[thin] (0.8,0) rectangle (6.4,3.8);

    % curva de aprendizado com os pontos observados
    \draw[thick] (0.8,1.0) -- (1.5,2.2) -- (2.4,2.9) -- (3.4,3.3)
                 -- (4.6,3.55) -- (6.4,3.7);
    \foreach \p in {(0.8,1.0),(1.5,2.2),(2.4,2.9),(3.4,3.3),(4.6,3.55),(6.4,3.7)}
        \fill \p circle (1.6pt);

    % linha do baseline (rótulo à direita, longe do ponteiro da AUC_ideal)
    \draw[dashed] (0,3.8) -- (7.0,3.8);
    \node[rot, anchor=south east] at (6.95,3.86)
        {\mbox{$\mathrm{Perf}_{\mathrm{baseline}}$} (supervisão completa)};

    % eixos
    \draw[seta] (-0.15,0) -- (7.35,0);
    \draw[seta] (0,-0.15) -- (0,4.6);
    \node[rot, anchor=south] at (0.3,4.6) {desempenho};
    \node[rot, anchor=north] at (3.6,-0.62) {rótulos consumidos ($|L_t|$)};
    \draw (0.8,0) -- (0.8,-0.12);
    \draw (6.4,0) -- (6.4,-0.12);
    \node[rot, anchor=north] at (0.8,-0.2) {$L_{\mathrm{ideal},0}$};
    \node[rot, anchor=north] at (6.4,-0.2) {$L_{\mathrm{final}}$};

    % rótulos das áreas: a fatia clara é estreita demais para texto dentro,
    % então a AUC_ideal é apontada de fora (achado da inspeção visual)
    \node[rot] at (4.5,1.7) {$\mathrm{AUC}_{\mathrm{real}}$};
    \node[rot, anchor=west] (auci) at (0.7,4.32)
        {$\mathrm{AUC}_{\mathrm{ideal}}$: o retângulo inteiro};
    \draw[seta, dashed] (1.6,4.15) -- (1.4,3.5);

    % ---------------- painel direito: fórmula e leitura --------------------
    \node[caixa, text width=58mm] (formula) at (10.6,3.5)
        {$\displaystyle \mathrm{LCE}=\frac{\mathrm{AUC}_{\mathrm{real}}}{\mathrm{AUC}_{\mathrm{ideal}}}$\\[3pt]
         \mbox{$\mathrm{AUC}_{\mathrm{ideal}}=\mathrm{Perf}_{\mathrm{baseline}}\times(L_{\mathrm{final}}-L_{\mathrm{ideal},0})$}\\[2pt]
         (Equação~\ref{eq:metodo-lce}; Apêndice~\ref{ap:lce})};
    \node[rot, text width=52mm, anchor=north, align=left] at (10.6,2.2)
        {$\mathrm{LCE}\approx 1$: a estratégia aproxima-se\\
         rapidamente do desempenho de\\
         referência com poucos rótulos.\\
         Integração numérica: regra de Simpson\\
         sobre os pontos observados (marcados).};
\end{tikzpicture}
\end{document}
%   (no modo standalone, troque as \ref{...} por texto fixo ou defina-as)
% Requer apenas arrows.meta e positioning (já em packages.tex). Legível em
% P&B: área sob a curva real = cinza escuro; retângulo ideal = cinza claro;
% pontos observados marcados (é sobre eles que Simpson integra).
% ============================================================================
\begin{tikzpicture}[
    x=1cm, y=1cm,
    caixa/.style={draw, rounded corners=2pt, align=center, text width=42mm,
                  minimum height=11mm, inner sep=3pt, font=\footnotesize},
    seta/.style={-{Stealth[length=2.4mm]}, thick},
    rot/.style={font=\scriptsize, align=center, inner sep=2pt}
]
    % ---------------- painel esquerdo: a geometria -------------------------
    \node[rot, anchor=west] at (-0.4,5.2)
        {\textbf{A geometria da LCE} (curva ilustrativa, sem dados medidos)};

    % áreas (desenhadas antes dos eixos)
    \fill[gray!8]  (0.8,0) rectangle (6.4,3.8);
    \fill[gray!25] (0.8,1.0) -- (1.5,2.2) -- (2.4,2.9) -- (3.4,3.3)
                   -- (4.6,3.55) -- (6.4,3.7) -- (6.4,0) -- (0.8,0) -- cycle;
    \draw[thin] (0.8,0) rectangle (6.4,3.8);

    % curva de aprendizado com os pontos observados
    \draw[thick] (0.8,1.0) -- (1.5,2.2) -- (2.4,2.9) -- (3.4,3.3)
                 -- (4.6,3.55) -- (6.4,3.7);
    \foreach \p in {(0.8,1.0),(1.5,2.2),(2.4,2.9),(3.4,3.3),(4.6,3.55),(6.4,3.7)}
        \fill \p circle (1.6pt);

    % linha do baseline (rótulo à direita, longe do ponteiro da AUC_ideal)
    \draw[dashed] (0,3.8) -- (7.0,3.8);
    \node[rot, anchor=south east] at (6.95,3.86)
        {\mbox{$\mathrm{Perf}_{\mathrm{baseline}}$} (supervisão completa)};

    % eixos
    \draw[seta] (-0.15,0) -- (7.35,0);
    \draw[seta] (0,-0.15) -- (0,4.6);
    \node[rot, anchor=south] at (0.3,4.6) {desempenho};
    \node[rot, anchor=north] at (3.6,-0.62) {rótulos consumidos ($|L_t|$)};
    \draw (0.8,0) -- (0.8,-0.12);
    \draw (6.4,0) -- (6.4,-0.12);
    \node[rot, anchor=north] at (0.8,-0.2) {$L_{\mathrm{ideal},0}$};
    \node[rot, anchor=north] at (6.4,-0.2) {$L_{\mathrm{final}}$};

    % rótulos das áreas: a fatia clara é estreita demais para texto dentro,
    % então a AUC_ideal é apontada de fora (achado da inspeção visual)
    \node[rot] at (4.5,1.7) {$\mathrm{AUC}_{\mathrm{real}}$};
    \node[rot, anchor=west] (auci) at (0.7,4.32)
        {$\mathrm{AUC}_{\mathrm{ideal}}$: o retângulo inteiro};
    \draw[seta, dashed] (1.6,4.15) -- (1.4,3.5);

    % ---------------- painel direito: fórmula e leitura --------------------
    \node[caixa, text width=58mm] (formula) at (10.6,3.5)
        {$\displaystyle \mathrm{LCE}=\frac{\mathrm{AUC}_{\mathrm{real}}}{\mathrm{AUC}_{\mathrm{ideal}}}$\\[3pt]
         \mbox{$\mathrm{AUC}_{\mathrm{ideal}}=\mathrm{Perf}_{\mathrm{baseline}}\times(L_{\mathrm{final}}-L_{\mathrm{ideal},0})$}\\[2pt]
         (Equação~\ref{eq:metodo-lce}; Apêndice~\ref{ap:lce})};
    \node[rot, text width=52mm, anchor=north, align=left] at (10.6,2.2)
        {$\mathrm{LCE}\approx 1$: a estratégia aproxima-se\\
         rapidamente do desempenho de\\
         referência com poucos rótulos.\\
         Integração numérica: regra de Simpson\\
         sobre os pontos observados (marcados).};
\end{tikzpicture}
\end{document}
%   (no modo standalone, troque as \ref{...} por texto fixo ou defina-as)
% Requer apenas arrows.meta e positioning (já em packages.tex). Legível em
% P&B: área sob a curva real = cinza escuro; retângulo ideal = cinza claro;
% pontos observados marcados (é sobre eles que Simpson integra).
% ============================================================================
\begin{tikzpicture}[
    x=1cm, y=1cm,
    caixa/.style={draw, rounded corners=2pt, align=center, text width=42mm,
                  minimum height=11mm, inner sep=3pt, font=\footnotesize},
    seta/.style={-{Stealth[length=2.4mm]}, thick},
    rot/.style={font=\scriptsize, align=center, inner sep=2pt}
]
    % ---------------- painel esquerdo: a geometria -------------------------
    \node[rot, anchor=west] at (-0.4,5.2)
        {\textbf{A geometria da LCE} (curva ilustrativa, sem dados medidos)};

    % áreas (desenhadas antes dos eixos)
    \fill[gray!8]  (0.8,0) rectangle (6.4,3.8);
    \fill[gray!25] (0.8,1.0) -- (1.5,2.2) -- (2.4,2.9) -- (3.4,3.3)
                   -- (4.6,3.55) -- (6.4,3.7) -- (6.4,0) -- (0.8,0) -- cycle;
    \draw[thin] (0.8,0) rectangle (6.4,3.8);

    % curva de aprendizado com os pontos observados
    \draw[thick] (0.8,1.0) -- (1.5,2.2) -- (2.4,2.9) -- (3.4,3.3)
                 -- (4.6,3.55) -- (6.4,3.7);
    \foreach \p in {(0.8,1.0),(1.5,2.2),(2.4,2.9),(3.4,3.3),(4.6,3.55),(6.4,3.7)}
        \fill \p circle (1.6pt);

    % linha do baseline (rótulo à direita, longe do ponteiro da AUC_ideal)
    \draw[dashed] (0,3.8) -- (7.0,3.8);
    \node[rot, anchor=south east] at (6.95,3.86)
        {\mbox{$\mathrm{Perf}_{\mathrm{baseline}}$} (supervisão completa)};

    % eixos
    \draw[seta] (-0.15,0) -- (7.35,0);
    \draw[seta] (0,-0.15) -- (0,4.6);
    \node[rot, anchor=south] at (0.3,4.6) {desempenho};
    \node[rot, anchor=north] at (3.6,-0.62) {rótulos consumidos ($|L_t|$)};
    \draw (0.8,0) -- (0.8,-0.12);
    \draw (6.4,0) -- (6.4,-0.12);
    \node[rot, anchor=north] at (0.8,-0.2) {$L_{\mathrm{ideal},0}$};
    \node[rot, anchor=north] at (6.4,-0.2) {$L_{\mathrm{final}}$};

    % rótulos das áreas: a fatia clara é estreita demais para texto dentro,
    % então a AUC_ideal é apontada de fora (achado da inspeção visual)
    \node[rot] at (4.5,1.7) {$\mathrm{AUC}_{\mathrm{real}}$};
    \node[rot, anchor=west] (auci) at (0.7,4.32)
        {$\mathrm{AUC}_{\mathrm{ideal}}$: o retângulo inteiro};
    \draw[seta, dashed] (1.6,4.15) -- (1.4,3.5);

    % ---------------- painel direito: fórmula e leitura --------------------
    \node[caixa, text width=58mm] (formula) at (10.6,3.5)
        {$\displaystyle \mathrm{LCE}=\frac{\mathrm{AUC}_{\mathrm{real}}}{\mathrm{AUC}_{\mathrm{ideal}}}$\\[3pt]
         \mbox{$\mathrm{AUC}_{\mathrm{ideal}}=\mathrm{Perf}_{\mathrm{baseline}}\times(L_{\mathrm{final}}-L_{\mathrm{ideal},0})$}\\[2pt]
         (Equação~\ref{eq:metodo-lce}; Apêndice~\ref{ap:lce})};
    \node[rot, text width=52mm, anchor=north, align=left] at (10.6,2.2)
        {$\mathrm{LCE}\approx 1$: a estratégia aproxima-se\\
         rapidamente do desempenho de\\
         referência com poucos rótulos.\\
         Integração numérica: regra de Simpson\\
         sobre os pontos observados (marcados).};
\end{tikzpicture}

%   A citação Guyon2011ALC (precedente ALC) permanece na prosa/legenda.
% Uso standalone (pré-visualização):
%   \documentclass[tikz,border=6pt]{standalone}
%   \usetikzlibrary{arrows.meta, positioning}
%   \begin{document}% ============================================================================
% ESQUEMA PROPOSTO (banca, sugestão — não mergeado): a geometria da LCE
% (Seção sec:metodo-metricas). A curva é ILUSTRATIVA e o desenho diz isso
% no título interno: nenhum dado medido, nenhum número além dos símbolos da
% própria equação — é a definição da métrica desenhada, não um resultado
% (compatível com o princípio V). Por decisão do autor (2026-08-23), a
% análise estatística e a precedência da ALC ficam SÓ na prosa.
%
% Uso no Cap. 3 (dentro de um ambiente figure):
%   % ============================================================================
% ESQUEMA PROPOSTO (banca, sugestão — não mergeado): a geometria da LCE
% (Seção sec:metodo-metricas). A curva é ILUSTRATIVA e o desenho diz isso
% no título interno: nenhum dado medido, nenhum número além dos símbolos da
% própria equação — é a definição da métrica desenhada, não um resultado
% (compatível com o princípio V). Por decisão do autor (2026-08-23), a
% análise estatística e a precedência da ALC ficam SÓ na prosa.
%
% Uso no Cap. 3 (dentro de um ambiente figure):
%   % ============================================================================
% ESQUEMA PROPOSTO (banca, sugestão — não mergeado): a geometria da LCE
% (Seção sec:metodo-metricas). A curva é ILUSTRATIVA e o desenho diz isso
% no título interno: nenhum dado medido, nenhum número além dos símbolos da
% própria equação — é a definição da métrica desenhada, não um resultado
% (compatível com o princípio V). Por decisão do autor (2026-08-23), a
% análise estatística e a precedência da ALC ficam SÓ na prosa.
%
% Uso no Cap. 3 (dentro de um ambiente figure):
%   \input{3-metodo/esquemas-propostos/esq-lce}
%   A citação Guyon2011ALC (precedente ALC) permanece na prosa/legenda.
% Uso standalone (pré-visualização):
%   \documentclass[tikz,border=6pt]{standalone}
%   \usetikzlibrary{arrows.meta, positioning}
%   \begin{document}\input{esq-lce.tex}\end{document}
%   (no modo standalone, troque as \ref{...} por texto fixo ou defina-as)
% Requer apenas arrows.meta e positioning (já em packages.tex). Legível em
% P&B: área sob a curva real = cinza escuro; retângulo ideal = cinza claro;
% pontos observados marcados (é sobre eles que Simpson integra).
% ============================================================================
\begin{tikzpicture}[
    x=1cm, y=1cm,
    caixa/.style={draw, rounded corners=2pt, align=center, text width=42mm,
                  minimum height=11mm, inner sep=3pt, font=\footnotesize},
    seta/.style={-{Stealth[length=2.4mm]}, thick},
    rot/.style={font=\scriptsize, align=center, inner sep=2pt}
]
    % ---------------- painel esquerdo: a geometria -------------------------
    \node[rot, anchor=west] at (-0.4,5.2)
        {\textbf{A geometria da LCE} (curva ilustrativa, sem dados medidos)};

    % áreas (desenhadas antes dos eixos)
    \fill[gray!8]  (0.8,0) rectangle (6.4,3.8);
    \fill[gray!25] (0.8,1.0) -- (1.5,2.2) -- (2.4,2.9) -- (3.4,3.3)
                   -- (4.6,3.55) -- (6.4,3.7) -- (6.4,0) -- (0.8,0) -- cycle;
    \draw[thin] (0.8,0) rectangle (6.4,3.8);

    % curva de aprendizado com os pontos observados
    \draw[thick] (0.8,1.0) -- (1.5,2.2) -- (2.4,2.9) -- (3.4,3.3)
                 -- (4.6,3.55) -- (6.4,3.7);
    \foreach \p in {(0.8,1.0),(1.5,2.2),(2.4,2.9),(3.4,3.3),(4.6,3.55),(6.4,3.7)}
        \fill \p circle (1.6pt);

    % linha do baseline (rótulo à direita, longe do ponteiro da AUC_ideal)
    \draw[dashed] (0,3.8) -- (7.0,3.8);
    \node[rot, anchor=south east] at (6.95,3.86)
        {\mbox{$\mathrm{Perf}_{\mathrm{baseline}}$} (supervisão completa)};

    % eixos
    \draw[seta] (-0.15,0) -- (7.35,0);
    \draw[seta] (0,-0.15) -- (0,4.6);
    \node[rot, anchor=south] at (0.3,4.6) {desempenho};
    \node[rot, anchor=north] at (3.6,-0.62) {rótulos consumidos ($|L_t|$)};
    \draw (0.8,0) -- (0.8,-0.12);
    \draw (6.4,0) -- (6.4,-0.12);
    \node[rot, anchor=north] at (0.8,-0.2) {$L_{\mathrm{ideal},0}$};
    \node[rot, anchor=north] at (6.4,-0.2) {$L_{\mathrm{final}}$};

    % rótulos das áreas: a fatia clara é estreita demais para texto dentro,
    % então a AUC_ideal é apontada de fora (achado da inspeção visual)
    \node[rot] at (4.5,1.7) {$\mathrm{AUC}_{\mathrm{real}}$};
    \node[rot, anchor=west] (auci) at (0.7,4.32)
        {$\mathrm{AUC}_{\mathrm{ideal}}$: o retângulo inteiro};
    \draw[seta, dashed] (1.6,4.15) -- (1.4,3.5);

    % ---------------- painel direito: fórmula e leitura --------------------
    \node[caixa, text width=58mm] (formula) at (10.6,3.5)
        {$\displaystyle \mathrm{LCE}=\frac{\mathrm{AUC}_{\mathrm{real}}}{\mathrm{AUC}_{\mathrm{ideal}}}$\\[3pt]
         \mbox{$\mathrm{AUC}_{\mathrm{ideal}}=\mathrm{Perf}_{\mathrm{baseline}}\times(L_{\mathrm{final}}-L_{\mathrm{ideal},0})$}\\[2pt]
         (Equação~\ref{eq:metodo-lce}; Apêndice~\ref{ap:lce})};
    \node[rot, text width=52mm, anchor=north, align=left] at (10.6,2.2)
        {$\mathrm{LCE}\approx 1$: a estratégia aproxima-se\\
         rapidamente do desempenho de\\
         referência com poucos rótulos.\\
         Integração numérica: regra de Simpson\\
         sobre os pontos observados (marcados).};
\end{tikzpicture}

%   A citação Guyon2011ALC (precedente ALC) permanece na prosa/legenda.
% Uso standalone (pré-visualização):
%   \documentclass[tikz,border=6pt]{standalone}
%   \usetikzlibrary{arrows.meta, positioning}
%   \begin{document}% ============================================================================
% ESQUEMA PROPOSTO (banca, sugestão — não mergeado): a geometria da LCE
% (Seção sec:metodo-metricas). A curva é ILUSTRATIVA e o desenho diz isso
% no título interno: nenhum dado medido, nenhum número além dos símbolos da
% própria equação — é a definição da métrica desenhada, não um resultado
% (compatível com o princípio V). Por decisão do autor (2026-08-23), a
% análise estatística e a precedência da ALC ficam SÓ na prosa.
%
% Uso no Cap. 3 (dentro de um ambiente figure):
%   \input{3-metodo/esquemas-propostos/esq-lce}
%   A citação Guyon2011ALC (precedente ALC) permanece na prosa/legenda.
% Uso standalone (pré-visualização):
%   \documentclass[tikz,border=6pt]{standalone}
%   \usetikzlibrary{arrows.meta, positioning}
%   \begin{document}\input{esq-lce.tex}\end{document}
%   (no modo standalone, troque as \ref{...} por texto fixo ou defina-as)
% Requer apenas arrows.meta e positioning (já em packages.tex). Legível em
% P&B: área sob a curva real = cinza escuro; retângulo ideal = cinza claro;
% pontos observados marcados (é sobre eles que Simpson integra).
% ============================================================================
\begin{tikzpicture}[
    x=1cm, y=1cm,
    caixa/.style={draw, rounded corners=2pt, align=center, text width=42mm,
                  minimum height=11mm, inner sep=3pt, font=\footnotesize},
    seta/.style={-{Stealth[length=2.4mm]}, thick},
    rot/.style={font=\scriptsize, align=center, inner sep=2pt}
]
    % ---------------- painel esquerdo: a geometria -------------------------
    \node[rot, anchor=west] at (-0.4,5.2)
        {\textbf{A geometria da LCE} (curva ilustrativa, sem dados medidos)};

    % áreas (desenhadas antes dos eixos)
    \fill[gray!8]  (0.8,0) rectangle (6.4,3.8);
    \fill[gray!25] (0.8,1.0) -- (1.5,2.2) -- (2.4,2.9) -- (3.4,3.3)
                   -- (4.6,3.55) -- (6.4,3.7) -- (6.4,0) -- (0.8,0) -- cycle;
    \draw[thin] (0.8,0) rectangle (6.4,3.8);

    % curva de aprendizado com os pontos observados
    \draw[thick] (0.8,1.0) -- (1.5,2.2) -- (2.4,2.9) -- (3.4,3.3)
                 -- (4.6,3.55) -- (6.4,3.7);
    \foreach \p in {(0.8,1.0),(1.5,2.2),(2.4,2.9),(3.4,3.3),(4.6,3.55),(6.4,3.7)}
        \fill \p circle (1.6pt);

    % linha do baseline (rótulo à direita, longe do ponteiro da AUC_ideal)
    \draw[dashed] (0,3.8) -- (7.0,3.8);
    \node[rot, anchor=south east] at (6.95,3.86)
        {\mbox{$\mathrm{Perf}_{\mathrm{baseline}}$} (supervisão completa)};

    % eixos
    \draw[seta] (-0.15,0) -- (7.35,0);
    \draw[seta] (0,-0.15) -- (0,4.6);
    \node[rot, anchor=south] at (0.3,4.6) {desempenho};
    \node[rot, anchor=north] at (3.6,-0.62) {rótulos consumidos ($|L_t|$)};
    \draw (0.8,0) -- (0.8,-0.12);
    \draw (6.4,0) -- (6.4,-0.12);
    \node[rot, anchor=north] at (0.8,-0.2) {$L_{\mathrm{ideal},0}$};
    \node[rot, anchor=north] at (6.4,-0.2) {$L_{\mathrm{final}}$};

    % rótulos das áreas: a fatia clara é estreita demais para texto dentro,
    % então a AUC_ideal é apontada de fora (achado da inspeção visual)
    \node[rot] at (4.5,1.7) {$\mathrm{AUC}_{\mathrm{real}}$};
    \node[rot, anchor=west] (auci) at (0.7,4.32)
        {$\mathrm{AUC}_{\mathrm{ideal}}$: o retângulo inteiro};
    \draw[seta, dashed] (1.6,4.15) -- (1.4,3.5);

    % ---------------- painel direito: fórmula e leitura --------------------
    \node[caixa, text width=58mm] (formula) at (10.6,3.5)
        {$\displaystyle \mathrm{LCE}=\frac{\mathrm{AUC}_{\mathrm{real}}}{\mathrm{AUC}_{\mathrm{ideal}}}$\\[3pt]
         \mbox{$\mathrm{AUC}_{\mathrm{ideal}}=\mathrm{Perf}_{\mathrm{baseline}}\times(L_{\mathrm{final}}-L_{\mathrm{ideal},0})$}\\[2pt]
         (Equação~\ref{eq:metodo-lce}; Apêndice~\ref{ap:lce})};
    \node[rot, text width=52mm, anchor=north, align=left] at (10.6,2.2)
        {$\mathrm{LCE}\approx 1$: a estratégia aproxima-se\\
         rapidamente do desempenho de\\
         referência com poucos rótulos.\\
         Integração numérica: regra de Simpson\\
         sobre os pontos observados (marcados).};
\end{tikzpicture}
\end{document}
%   (no modo standalone, troque as \ref{...} por texto fixo ou defina-as)
% Requer apenas arrows.meta e positioning (já em packages.tex). Legível em
% P&B: área sob a curva real = cinza escuro; retângulo ideal = cinza claro;
% pontos observados marcados (é sobre eles que Simpson integra).
% ============================================================================
\begin{tikzpicture}[
    x=1cm, y=1cm,
    caixa/.style={draw, rounded corners=2pt, align=center, text width=42mm,
                  minimum height=11mm, inner sep=3pt, font=\footnotesize},
    seta/.style={-{Stealth[length=2.4mm]}, thick},
    rot/.style={font=\scriptsize, align=center, inner sep=2pt}
]
    % ---------------- painel esquerdo: a geometria -------------------------
    \node[rot, anchor=west] at (-0.4,5.2)
        {\textbf{A geometria da LCE} (curva ilustrativa, sem dados medidos)};

    % áreas (desenhadas antes dos eixos)
    \fill[gray!8]  (0.8,0) rectangle (6.4,3.8);
    \fill[gray!25] (0.8,1.0) -- (1.5,2.2) -- (2.4,2.9) -- (3.4,3.3)
                   -- (4.6,3.55) -- (6.4,3.7) -- (6.4,0) -- (0.8,0) -- cycle;
    \draw[thin] (0.8,0) rectangle (6.4,3.8);

    % curva de aprendizado com os pontos observados
    \draw[thick] (0.8,1.0) -- (1.5,2.2) -- (2.4,2.9) -- (3.4,3.3)
                 -- (4.6,3.55) -- (6.4,3.7);
    \foreach \p in {(0.8,1.0),(1.5,2.2),(2.4,2.9),(3.4,3.3),(4.6,3.55),(6.4,3.7)}
        \fill \p circle (1.6pt);

    % linha do baseline (rótulo à direita, longe do ponteiro da AUC_ideal)
    \draw[dashed] (0,3.8) -- (7.0,3.8);
    \node[rot, anchor=south east] at (6.95,3.86)
        {\mbox{$\mathrm{Perf}_{\mathrm{baseline}}$} (supervisão completa)};

    % eixos
    \draw[seta] (-0.15,0) -- (7.35,0);
    \draw[seta] (0,-0.15) -- (0,4.6);
    \node[rot, anchor=south] at (0.3,4.6) {desempenho};
    \node[rot, anchor=north] at (3.6,-0.62) {rótulos consumidos ($|L_t|$)};
    \draw (0.8,0) -- (0.8,-0.12);
    \draw (6.4,0) -- (6.4,-0.12);
    \node[rot, anchor=north] at (0.8,-0.2) {$L_{\mathrm{ideal},0}$};
    \node[rot, anchor=north] at (6.4,-0.2) {$L_{\mathrm{final}}$};

    % rótulos das áreas: a fatia clara é estreita demais para texto dentro,
    % então a AUC_ideal é apontada de fora (achado da inspeção visual)
    \node[rot] at (4.5,1.7) {$\mathrm{AUC}_{\mathrm{real}}$};
    \node[rot, anchor=west] (auci) at (0.7,4.32)
        {$\mathrm{AUC}_{\mathrm{ideal}}$: o retângulo inteiro};
    \draw[seta, dashed] (1.6,4.15) -- (1.4,3.5);

    % ---------------- painel direito: fórmula e leitura --------------------
    \node[caixa, text width=58mm] (formula) at (10.6,3.5)
        {$\displaystyle \mathrm{LCE}=\frac{\mathrm{AUC}_{\mathrm{real}}}{\mathrm{AUC}_{\mathrm{ideal}}}$\\[3pt]
         \mbox{$\mathrm{AUC}_{\mathrm{ideal}}=\mathrm{Perf}_{\mathrm{baseline}}\times(L_{\mathrm{final}}-L_{\mathrm{ideal},0})$}\\[2pt]
         (Equação~\ref{eq:metodo-lce}; Apêndice~\ref{ap:lce})};
    \node[rot, text width=52mm, anchor=north, align=left] at (10.6,2.2)
        {$\mathrm{LCE}\approx 1$: a estratégia aproxima-se\\
         rapidamente do desempenho de\\
         referência com poucos rótulos.\\
         Integração numérica: regra de Simpson\\
         sobre os pontos observados (marcados).};
\end{tikzpicture}

%   A citação Guyon2011ALC (precedente ALC) permanece na prosa/legenda.
% Uso standalone (pré-visualização):
%   \documentclass[tikz,border=6pt]{standalone}
%   \usetikzlibrary{arrows.meta, positioning}
%   \begin{document}% ============================================================================
% ESQUEMA PROPOSTO (banca, sugestão — não mergeado): a geometria da LCE
% (Seção sec:metodo-metricas). A curva é ILUSTRATIVA e o desenho diz isso
% no título interno: nenhum dado medido, nenhum número além dos símbolos da
% própria equação — é a definição da métrica desenhada, não um resultado
% (compatível com o princípio V). Por decisão do autor (2026-08-23), a
% análise estatística e a precedência da ALC ficam SÓ na prosa.
%
% Uso no Cap. 3 (dentro de um ambiente figure):
%   % ============================================================================
% ESQUEMA PROPOSTO (banca, sugestão — não mergeado): a geometria da LCE
% (Seção sec:metodo-metricas). A curva é ILUSTRATIVA e o desenho diz isso
% no título interno: nenhum dado medido, nenhum número além dos símbolos da
% própria equação — é a definição da métrica desenhada, não um resultado
% (compatível com o princípio V). Por decisão do autor (2026-08-23), a
% análise estatística e a precedência da ALC ficam SÓ na prosa.
%
% Uso no Cap. 3 (dentro de um ambiente figure):
%   \input{3-metodo/esquemas-propostos/esq-lce}
%   A citação Guyon2011ALC (precedente ALC) permanece na prosa/legenda.
% Uso standalone (pré-visualização):
%   \documentclass[tikz,border=6pt]{standalone}
%   \usetikzlibrary{arrows.meta, positioning}
%   \begin{document}\input{esq-lce.tex}\end{document}
%   (no modo standalone, troque as \ref{...} por texto fixo ou defina-as)
% Requer apenas arrows.meta e positioning (já em packages.tex). Legível em
% P&B: área sob a curva real = cinza escuro; retângulo ideal = cinza claro;
% pontos observados marcados (é sobre eles que Simpson integra).
% ============================================================================
\begin{tikzpicture}[
    x=1cm, y=1cm,
    caixa/.style={draw, rounded corners=2pt, align=center, text width=42mm,
                  minimum height=11mm, inner sep=3pt, font=\footnotesize},
    seta/.style={-{Stealth[length=2.4mm]}, thick},
    rot/.style={font=\scriptsize, align=center, inner sep=2pt}
]
    % ---------------- painel esquerdo: a geometria -------------------------
    \node[rot, anchor=west] at (-0.4,5.2)
        {\textbf{A geometria da LCE} (curva ilustrativa, sem dados medidos)};

    % áreas (desenhadas antes dos eixos)
    \fill[gray!8]  (0.8,0) rectangle (6.4,3.8);
    \fill[gray!25] (0.8,1.0) -- (1.5,2.2) -- (2.4,2.9) -- (3.4,3.3)
                   -- (4.6,3.55) -- (6.4,3.7) -- (6.4,0) -- (0.8,0) -- cycle;
    \draw[thin] (0.8,0) rectangle (6.4,3.8);

    % curva de aprendizado com os pontos observados
    \draw[thick] (0.8,1.0) -- (1.5,2.2) -- (2.4,2.9) -- (3.4,3.3)
                 -- (4.6,3.55) -- (6.4,3.7);
    \foreach \p in {(0.8,1.0),(1.5,2.2),(2.4,2.9),(3.4,3.3),(4.6,3.55),(6.4,3.7)}
        \fill \p circle (1.6pt);

    % linha do baseline (rótulo à direita, longe do ponteiro da AUC_ideal)
    \draw[dashed] (0,3.8) -- (7.0,3.8);
    \node[rot, anchor=south east] at (6.95,3.86)
        {\mbox{$\mathrm{Perf}_{\mathrm{baseline}}$} (supervisão completa)};

    % eixos
    \draw[seta] (-0.15,0) -- (7.35,0);
    \draw[seta] (0,-0.15) -- (0,4.6);
    \node[rot, anchor=south] at (0.3,4.6) {desempenho};
    \node[rot, anchor=north] at (3.6,-0.62) {rótulos consumidos ($|L_t|$)};
    \draw (0.8,0) -- (0.8,-0.12);
    \draw (6.4,0) -- (6.4,-0.12);
    \node[rot, anchor=north] at (0.8,-0.2) {$L_{\mathrm{ideal},0}$};
    \node[rot, anchor=north] at (6.4,-0.2) {$L_{\mathrm{final}}$};

    % rótulos das áreas: a fatia clara é estreita demais para texto dentro,
    % então a AUC_ideal é apontada de fora (achado da inspeção visual)
    \node[rot] at (4.5,1.7) {$\mathrm{AUC}_{\mathrm{real}}$};
    \node[rot, anchor=west] (auci) at (0.7,4.32)
        {$\mathrm{AUC}_{\mathrm{ideal}}$: o retângulo inteiro};
    \draw[seta, dashed] (1.6,4.15) -- (1.4,3.5);

    % ---------------- painel direito: fórmula e leitura --------------------
    \node[caixa, text width=58mm] (formula) at (10.6,3.5)
        {$\displaystyle \mathrm{LCE}=\frac{\mathrm{AUC}_{\mathrm{real}}}{\mathrm{AUC}_{\mathrm{ideal}}}$\\[3pt]
         \mbox{$\mathrm{AUC}_{\mathrm{ideal}}=\mathrm{Perf}_{\mathrm{baseline}}\times(L_{\mathrm{final}}-L_{\mathrm{ideal},0})$}\\[2pt]
         (Equação~\ref{eq:metodo-lce}; Apêndice~\ref{ap:lce})};
    \node[rot, text width=52mm, anchor=north, align=left] at (10.6,2.2)
        {$\mathrm{LCE}\approx 1$: a estratégia aproxima-se\\
         rapidamente do desempenho de\\
         referência com poucos rótulos.\\
         Integração numérica: regra de Simpson\\
         sobre os pontos observados (marcados).};
\end{tikzpicture}

%   A citação Guyon2011ALC (precedente ALC) permanece na prosa/legenda.
% Uso standalone (pré-visualização):
%   \documentclass[tikz,border=6pt]{standalone}
%   \usetikzlibrary{arrows.meta, positioning}
%   \begin{document}% ============================================================================
% ESQUEMA PROPOSTO (banca, sugestão — não mergeado): a geometria da LCE
% (Seção sec:metodo-metricas). A curva é ILUSTRATIVA e o desenho diz isso
% no título interno: nenhum dado medido, nenhum número além dos símbolos da
% própria equação — é a definição da métrica desenhada, não um resultado
% (compatível com o princípio V). Por decisão do autor (2026-08-23), a
% análise estatística e a precedência da ALC ficam SÓ na prosa.
%
% Uso no Cap. 3 (dentro de um ambiente figure):
%   \input{3-metodo/esquemas-propostos/esq-lce}
%   A citação Guyon2011ALC (precedente ALC) permanece na prosa/legenda.
% Uso standalone (pré-visualização):
%   \documentclass[tikz,border=6pt]{standalone}
%   \usetikzlibrary{arrows.meta, positioning}
%   \begin{document}\input{esq-lce.tex}\end{document}
%   (no modo standalone, troque as \ref{...} por texto fixo ou defina-as)
% Requer apenas arrows.meta e positioning (já em packages.tex). Legível em
% P&B: área sob a curva real = cinza escuro; retângulo ideal = cinza claro;
% pontos observados marcados (é sobre eles que Simpson integra).
% ============================================================================
\begin{tikzpicture}[
    x=1cm, y=1cm,
    caixa/.style={draw, rounded corners=2pt, align=center, text width=42mm,
                  minimum height=11mm, inner sep=3pt, font=\footnotesize},
    seta/.style={-{Stealth[length=2.4mm]}, thick},
    rot/.style={font=\scriptsize, align=center, inner sep=2pt}
]
    % ---------------- painel esquerdo: a geometria -------------------------
    \node[rot, anchor=west] at (-0.4,5.2)
        {\textbf{A geometria da LCE} (curva ilustrativa, sem dados medidos)};

    % áreas (desenhadas antes dos eixos)
    \fill[gray!8]  (0.8,0) rectangle (6.4,3.8);
    \fill[gray!25] (0.8,1.0) -- (1.5,2.2) -- (2.4,2.9) -- (3.4,3.3)
                   -- (4.6,3.55) -- (6.4,3.7) -- (6.4,0) -- (0.8,0) -- cycle;
    \draw[thin] (0.8,0) rectangle (6.4,3.8);

    % curva de aprendizado com os pontos observados
    \draw[thick] (0.8,1.0) -- (1.5,2.2) -- (2.4,2.9) -- (3.4,3.3)
                 -- (4.6,3.55) -- (6.4,3.7);
    \foreach \p in {(0.8,1.0),(1.5,2.2),(2.4,2.9),(3.4,3.3),(4.6,3.55),(6.4,3.7)}
        \fill \p circle (1.6pt);

    % linha do baseline (rótulo à direita, longe do ponteiro da AUC_ideal)
    \draw[dashed] (0,3.8) -- (7.0,3.8);
    \node[rot, anchor=south east] at (6.95,3.86)
        {\mbox{$\mathrm{Perf}_{\mathrm{baseline}}$} (supervisão completa)};

    % eixos
    \draw[seta] (-0.15,0) -- (7.35,0);
    \draw[seta] (0,-0.15) -- (0,4.6);
    \node[rot, anchor=south] at (0.3,4.6) {desempenho};
    \node[rot, anchor=north] at (3.6,-0.62) {rótulos consumidos ($|L_t|$)};
    \draw (0.8,0) -- (0.8,-0.12);
    \draw (6.4,0) -- (6.4,-0.12);
    \node[rot, anchor=north] at (0.8,-0.2) {$L_{\mathrm{ideal},0}$};
    \node[rot, anchor=north] at (6.4,-0.2) {$L_{\mathrm{final}}$};

    % rótulos das áreas: a fatia clara é estreita demais para texto dentro,
    % então a AUC_ideal é apontada de fora (achado da inspeção visual)
    \node[rot] at (4.5,1.7) {$\mathrm{AUC}_{\mathrm{real}}$};
    \node[rot, anchor=west] (auci) at (0.7,4.32)
        {$\mathrm{AUC}_{\mathrm{ideal}}$: o retângulo inteiro};
    \draw[seta, dashed] (1.6,4.15) -- (1.4,3.5);

    % ---------------- painel direito: fórmula e leitura --------------------
    \node[caixa, text width=58mm] (formula) at (10.6,3.5)
        {$\displaystyle \mathrm{LCE}=\frac{\mathrm{AUC}_{\mathrm{real}}}{\mathrm{AUC}_{\mathrm{ideal}}}$\\[3pt]
         \mbox{$\mathrm{AUC}_{\mathrm{ideal}}=\mathrm{Perf}_{\mathrm{baseline}}\times(L_{\mathrm{final}}-L_{\mathrm{ideal},0})$}\\[2pt]
         (Equação~\ref{eq:metodo-lce}; Apêndice~\ref{ap:lce})};
    \node[rot, text width=52mm, anchor=north, align=left] at (10.6,2.2)
        {$\mathrm{LCE}\approx 1$: a estratégia aproxima-se\\
         rapidamente do desempenho de\\
         referência com poucos rótulos.\\
         Integração numérica: regra de Simpson\\
         sobre os pontos observados (marcados).};
\end{tikzpicture}
\end{document}
%   (no modo standalone, troque as \ref{...} por texto fixo ou defina-as)
% Requer apenas arrows.meta e positioning (já em packages.tex). Legível em
% P&B: área sob a curva real = cinza escuro; retângulo ideal = cinza claro;
% pontos observados marcados (é sobre eles que Simpson integra).
% ============================================================================
\begin{tikzpicture}[
    x=1cm, y=1cm,
    caixa/.style={draw, rounded corners=2pt, align=center, text width=42mm,
                  minimum height=11mm, inner sep=3pt, font=\footnotesize},
    seta/.style={-{Stealth[length=2.4mm]}, thick},
    rot/.style={font=\scriptsize, align=center, inner sep=2pt}
]
    % ---------------- painel esquerdo: a geometria -------------------------
    \node[rot, anchor=west] at (-0.4,5.2)
        {\textbf{A geometria da LCE} (curva ilustrativa, sem dados medidos)};

    % áreas (desenhadas antes dos eixos)
    \fill[gray!8]  (0.8,0) rectangle (6.4,3.8);
    \fill[gray!25] (0.8,1.0) -- (1.5,2.2) -- (2.4,2.9) -- (3.4,3.3)
                   -- (4.6,3.55) -- (6.4,3.7) -- (6.4,0) -- (0.8,0) -- cycle;
    \draw[thin] (0.8,0) rectangle (6.4,3.8);

    % curva de aprendizado com os pontos observados
    \draw[thick] (0.8,1.0) -- (1.5,2.2) -- (2.4,2.9) -- (3.4,3.3)
                 -- (4.6,3.55) -- (6.4,3.7);
    \foreach \p in {(0.8,1.0),(1.5,2.2),(2.4,2.9),(3.4,3.3),(4.6,3.55),(6.4,3.7)}
        \fill \p circle (1.6pt);

    % linha do baseline (rótulo à direita, longe do ponteiro da AUC_ideal)
    \draw[dashed] (0,3.8) -- (7.0,3.8);
    \node[rot, anchor=south east] at (6.95,3.86)
        {\mbox{$\mathrm{Perf}_{\mathrm{baseline}}$} (supervisão completa)};

    % eixos
    \draw[seta] (-0.15,0) -- (7.35,0);
    \draw[seta] (0,-0.15) -- (0,4.6);
    \node[rot, anchor=south] at (0.3,4.6) {desempenho};
    \node[rot, anchor=north] at (3.6,-0.62) {rótulos consumidos ($|L_t|$)};
    \draw (0.8,0) -- (0.8,-0.12);
    \draw (6.4,0) -- (6.4,-0.12);
    \node[rot, anchor=north] at (0.8,-0.2) {$L_{\mathrm{ideal},0}$};
    \node[rot, anchor=north] at (6.4,-0.2) {$L_{\mathrm{final}}$};

    % rótulos das áreas: a fatia clara é estreita demais para texto dentro,
    % então a AUC_ideal é apontada de fora (achado da inspeção visual)
    \node[rot] at (4.5,1.7) {$\mathrm{AUC}_{\mathrm{real}}$};
    \node[rot, anchor=west] (auci) at (0.7,4.32)
        {$\mathrm{AUC}_{\mathrm{ideal}}$: o retângulo inteiro};
    \draw[seta, dashed] (1.6,4.15) -- (1.4,3.5);

    % ---------------- painel direito: fórmula e leitura --------------------
    \node[caixa, text width=58mm] (formula) at (10.6,3.5)
        {$\displaystyle \mathrm{LCE}=\frac{\mathrm{AUC}_{\mathrm{real}}}{\mathrm{AUC}_{\mathrm{ideal}}}$\\[3pt]
         \mbox{$\mathrm{AUC}_{\mathrm{ideal}}=\mathrm{Perf}_{\mathrm{baseline}}\times(L_{\mathrm{final}}-L_{\mathrm{ideal},0})$}\\[2pt]
         (Equação~\ref{eq:metodo-lce}; Apêndice~\ref{ap:lce})};
    \node[rot, text width=52mm, anchor=north, align=left] at (10.6,2.2)
        {$\mathrm{LCE}\approx 1$: a estratégia aproxima-se\\
         rapidamente do desempenho de\\
         referência com poucos rótulos.\\
         Integração numérica: regra de Simpson\\
         sobre os pontos observados (marcados).};
\end{tikzpicture}
\end{document}
%   (no modo standalone, troque as \ref{...} por texto fixo ou defina-as)
% Requer apenas arrows.meta e positioning (já em packages.tex). Legível em
% P&B: área sob a curva real = cinza escuro; retângulo ideal = cinza claro;
% pontos observados marcados (é sobre eles que Simpson integra).
% ============================================================================
\begin{tikzpicture}[
    x=1cm, y=1cm,
    caixa/.style={draw, rounded corners=2pt, align=center, text width=42mm,
                  minimum height=11mm, inner sep=3pt, font=\footnotesize},
    seta/.style={-{Stealth[length=2.4mm]}, thick},
    rot/.style={font=\scriptsize, align=center, inner sep=2pt}
]
    % ---------------- painel esquerdo: a geometria -------------------------
    \node[rot, anchor=west] at (-0.4,5.2)
        {\textbf{A geometria da LCE} (curva ilustrativa, sem dados medidos)};

    % áreas (desenhadas antes dos eixos)
    \fill[gray!8]  (0.8,0) rectangle (6.4,3.8);
    \fill[gray!25] (0.8,1.0) -- (1.5,2.2) -- (2.4,2.9) -- (3.4,3.3)
                   -- (4.6,3.55) -- (6.4,3.7) -- (6.4,0) -- (0.8,0) -- cycle;
    \draw[thin] (0.8,0) rectangle (6.4,3.8);

    % curva de aprendizado com os pontos observados
    \draw[thick] (0.8,1.0) -- (1.5,2.2) -- (2.4,2.9) -- (3.4,3.3)
                 -- (4.6,3.55) -- (6.4,3.7);
    \foreach \p in {(0.8,1.0),(1.5,2.2),(2.4,2.9),(3.4,3.3),(4.6,3.55),(6.4,3.7)}
        \fill \p circle (1.6pt);

    % linha do baseline (rótulo à direita, longe do ponteiro da AUC_ideal)
    \draw[dashed] (0,3.8) -- (7.0,3.8);
    \node[rot, anchor=south east] at (6.95,3.86)
        {\mbox{$\mathrm{Perf}_{\mathrm{baseline}}$} (supervisão completa)};

    % eixos
    \draw[seta] (-0.15,0) -- (7.35,0);
    \draw[seta] (0,-0.15) -- (0,4.6);
    \node[rot, anchor=south] at (0.3,4.6) {desempenho};
    \node[rot, anchor=north] at (3.6,-0.62) {rótulos consumidos ($|L_t|$)};
    \draw (0.8,0) -- (0.8,-0.12);
    \draw (6.4,0) -- (6.4,-0.12);
    \node[rot, anchor=north] at (0.8,-0.2) {$L_{\mathrm{ideal},0}$};
    \node[rot, anchor=north] at (6.4,-0.2) {$L_{\mathrm{final}}$};

    % rótulos das áreas: a fatia clara é estreita demais para texto dentro,
    % então a AUC_ideal é apontada de fora (achado da inspeção visual)
    \node[rot] at (4.5,1.7) {$\mathrm{AUC}_{\mathrm{real}}$};
    \node[rot, anchor=west] (auci) at (0.7,4.32)
        {$\mathrm{AUC}_{\mathrm{ideal}}$: o retângulo inteiro};
    \draw[seta, dashed] (1.6,4.15) -- (1.4,3.5);

    % ---------------- painel direito: fórmula e leitura --------------------
    \node[caixa, text width=58mm] (formula) at (10.6,3.5)
        {$\displaystyle \mathrm{LCE}=\frac{\mathrm{AUC}_{\mathrm{real}}}{\mathrm{AUC}_{\mathrm{ideal}}}$\\[3pt]
         \mbox{$\mathrm{AUC}_{\mathrm{ideal}}=\mathrm{Perf}_{\mathrm{baseline}}\times(L_{\mathrm{final}}-L_{\mathrm{ideal},0})$}\\[2pt]
         (Equação~\ref{eq:metodo-lce}; Apêndice~\ref{ap:lce})};
    \node[rot, text width=52mm, anchor=north, align=left] at (10.6,2.2)
        {$\mathrm{LCE}\approx 1$: a estratégia aproxima-se\\
         rapidamente do desempenho de\\
         referência com poucos rótulos.\\
         Integração numérica: regra de Simpson\\
         sobre os pontos observados (marcados).};
\end{tikzpicture}
\end{document}
%   (no modo standalone, troque as \ref{...} por texto fixo ou defina-as)
% Requer apenas arrows.meta e positioning (já em packages.tex). Legível em
% P&B: área sob a curva real = cinza escuro; retângulo ideal = cinza claro;
% pontos observados marcados (é sobre eles que Simpson integra).
% ============================================================================
\begin{tikzpicture}[
    x=1cm, y=1cm,
    caixa/.style={draw, rounded corners=2pt, align=center, text width=42mm,
                  minimum height=11mm, inner sep=3pt, font=\footnotesize},
    seta/.style={-{Stealth[length=2.4mm]}, thick},
    rot/.style={font=\scriptsize, align=center, inner sep=2pt}
]
    % ---------------- painel esquerdo: a geometria -------------------------
    \node[rot, anchor=west] at (-0.4,5.2)
        {\textbf{A geometria da LCE} (curva ilustrativa, sem dados medidos)};

    % áreas (desenhadas antes dos eixos)
    \fill[gray!8]  (0.8,0) rectangle (6.4,3.8);
    \fill[gray!25] (0.8,1.0) -- (1.5,2.2) -- (2.4,2.9) -- (3.4,3.3)
                   -- (4.6,3.55) -- (6.4,3.7) -- (6.4,0) -- (0.8,0) -- cycle;
    \draw[thin] (0.8,0) rectangle (6.4,3.8);

    % curva de aprendizado com os pontos observados
    \draw[thick] (0.8,1.0) -- (1.5,2.2) -- (2.4,2.9) -- (3.4,3.3)
                 -- (4.6,3.55) -- (6.4,3.7);
    \foreach \p in {(0.8,1.0),(1.5,2.2),(2.4,2.9),(3.4,3.3),(4.6,3.55),(6.4,3.7)}
        \fill \p circle (1.6pt);

    % linha do baseline (rótulo à direita, longe do ponteiro da AUC_ideal)
    \draw[dashed] (0,3.8) -- (7.0,3.8);
    \node[rot, anchor=south east] at (6.95,3.86)
        {\mbox{$\mathrm{Perf}_{\mathrm{baseline}}$} (supervisão completa)};

    % eixos
    \draw[seta] (-0.15,0) -- (7.35,0);
    \draw[seta] (0,-0.15) -- (0,4.6);
    \node[rot, anchor=south] at (0.3,4.6) {desempenho};
    \node[rot, anchor=north] at (3.6,-0.62) {rótulos consumidos ($|L_t|$)};
    \draw (0.8,0) -- (0.8,-0.12);
    \draw (6.4,0) -- (6.4,-0.12);
    \node[rot, anchor=north] at (0.8,-0.2) {$L_{\mathrm{ideal},0}$};
    \node[rot, anchor=north] at (6.4,-0.2) {$L_{\mathrm{final}}$};

    % rótulos das áreas: a fatia clara é estreita demais para texto dentro,
    % então a AUC_ideal é apontada de fora (achado da inspeção visual)
    \node[rot] at (4.5,1.7) {$\mathrm{AUC}_{\mathrm{real}}$};
    \node[rot, anchor=west] (auci) at (0.7,4.32)
        {$\mathrm{AUC}_{\mathrm{ideal}}$: o retângulo inteiro};
    \draw[seta, dashed] (1.6,4.15) -- (1.4,3.5);

    % ---------------- painel direito: fórmula e leitura --------------------
    \node[caixa, text width=58mm] (formula) at (10.6,3.5)
        {$\displaystyle \mathrm{LCE}=\frac{\mathrm{AUC}_{\mathrm{real}}}{\mathrm{AUC}_{\mathrm{ideal}}}$\\[3pt]
         \mbox{$\mathrm{AUC}_{\mathrm{ideal}}=\mathrm{Perf}_{\mathrm{baseline}}\times(L_{\mathrm{final}}-L_{\mathrm{ideal},0})$}\\[2pt]
         (Equação~\ref{eq:metodo-lce}; Apêndice~\ref{ap:lce})};
    \node[rot, text width=52mm, anchor=north, align=left] at (10.6,2.2)
        {$\mathrm{LCE}\approx 1$: a estratégia aproxima-se\\
         rapidamente do desempenho de\\
         referência com poucos rótulos.\\
         Integração numérica: regra de Simpson\\
         sobre os pontos observados (marcados).};
\end{tikzpicture}
